Txosten honetan, Robotika eta Kontrol Adimendunak irakasgaiko sentsoreen praktikan garatu ditugun bi portaera autonomoak aurkezten ditugu, ROSMASTER X3 robotarentzat eta ROS~2 erabiliz.

Lehenengo portaerak oztopoak ekiditen ditu LiDARraren irakurketetan oinarrituta. Robotaren ingurua hiru sektoretan banatu dugu (aurrea, ezkerra eta eskuina), sektore bakoitzean oztoporik hurbilena neurtzen dugu, eta histeresia duen bi egoerako makina batek erabakitzen du aurrera egin ala lekuan bertan biratu. Gazeboko \texttt{simple} munduan, robotak \bete{X} minutu eman ditu talkarik egin gabe. Haren esplorazioa, ordea, \bete{eremu jakin batera mugatzen da}, portaera erreaktiboa, deterministikoa eta memoriarik gabekoa delako.

Bigarren portaerak errepideari jarraitzen dio, aurreko kameraren irudi segmentatua erabiliz. Irudiaren banda horizontal batean errepidearen zentroidea kalkulatzen dugu, eta robotak harantz biratzen du, errorearekiko proportzionala den kontrol batekin. Irabazia abiadura maximoarekiko proportzionala da, ibilbidea abiaduraren mende egon ez dadin. Simulazioan, robotak \bete{N} bira osatu ditu \bete{arazorik gabe}. Robot errealean, bost saiakeratatik \bete{N} izan dira arrakastatsuak (\bete{\%}), eta bira bakoitzak batez beste \bete{T}~s iraun du.
